\textbf{COLUNI 2021 -- Questão 28}\par\medskip A polinização biótica é considerada um serviço ambiental de grande importância para a biodiversidade vegetal. Nesse processo, espécies animais, especialmente abelhas, visitam flores em busca de recursos como néctar e pólen e, acidentalmente, tocam estruturas florais carreando assim grãos de pólen de uma flor para outra. Sobre a polinização e a biodiversidade vegetal, é INCORRETO afirmar:\begin{enumerate}
\item Espécies vegetais pertencentes ao grupo das Angiospermas caracterizam-se pela presença de flores e de
       frutos e representam grande avanço no processo evolutivo, dado o aumento da variabilidade genética
       proporcionado pela polinização cruzada.
\item As estruturas florais mencionadas são as anteras, que compõem a porção masculina da flor, nas quais
       são produzidos os grãos de pólen; e os filetes, que compõem a porção feminina da flor, nos quais estão
       localizados os nectários florais.
\item Existem espécies com diferentes arranjos florais, por exemplo, espécies com apenas flores hermafroditas,
       com flores masculinas e femininas na mesma planta ou com flores masculinas e femininas em plantas
       diferentes.
\item O fruto é uma estrutura importante no processo de dispersão das sementes, originando-se a partir do
       desenvolvimento do ovário, que abriga as sementes, formadas pela fecundação do óvulo pelo gameta
       masculino, presente no grão de pólen.\end{enumerate}
